\section{Behavioral targets and trajectory supervision}
\label{sec:C08:sft}
Post-training changes a pretrained conditional model to support a more specific behavior distribution. For agents, the target includes selecting tools, constructing valid arguments, using observations, recognizing failure, and stopping under appropriate conditions. Supervised fine-tuning minimizes negative log likelihood on selected target actions or assistant tokens. Its statistical meaning is imitation of the recorded conditional distribution. A successful demonstration supplies one feasible behavior, not proof that every action in it is necessary or optimal. If a trace succeeds despite an unnecessary retry, direct imitation can preserve that inefficiency.

Let a trajectory contain contexts $C_t$ and actions $a_t$. A sequence-level objective can weight each trajectory equally, while a token-level objective weights long outputs more heavily. These objectives differ when demonstrations have different reasoning lengths or numbers of tool calls. Report the weighting and loss mask, including whether tool observations contribute. A dataset that trains the model to predict authoritative-looking tool results can encourage simulation of results unless the output interface clearly distinguishes an action request from an observation. The model should not be rewarded for inventing the external evidence it is supposed to obtain.

Failure data require deliberate targets. A failed trajectory can teach recognition and correction without making its original erroneous actions desirable. One approach supervises only the corrective continuation given the failed prefix; another trains outcome or process evaluators. Simply concatenating failed and successful traces under the same imitation objective mixes incompatible behavioral targets. The data record should identify which action was erroneous, what feedback was available, and whether the correction depended on later information. Retrospective reasoning can be useful as an explanation while still being invalid as the conditioning prefix for the original action.

Instruction following through human feedback provides an influential post-training baseline, but its preference judgments are specific to the task and annotator protocol~\cite{ouyang2022instructgpt}. An agent policy needs additional evaluation of external effects and long-run state. A response preferred by a human reader can still call the wrong tool or omit a necessary verification step. Consequently, preference quality and executable correctness should be recorded separately. The most convincing training signal combines task-relevant judgment with evidence from the actual environment, while preserving uncertainty when the specification itself is ambiguous.

